\documentclass[addpoints,answers]{exam}

\usepackage[utf8]{inputenc}
\usepackage{calc}
\usepackage{listings}
\usepackage{xcolor}
\usepackage{xspace}

\definecolor{codegreen}{rgb}{0,0.6,0}
\definecolor{codegray}{rgb}{0.5,0.5,0.5}
\definecolor{codepurple}{rgb}{0.58,0,0.82}

\lstdefinestyle{Cstyle}{
    commentstyle=\color{codegreen},
    keywordstyle=\color{magenta},
    numberstyle=\tiny\color{codegray},
    stringstyle=\color{codepurple},
    basicstyle=\ttfamily\small,
    breakatwhitespace=true,         
    breaklines=true,                 
    captionpos=b,                    
    keepspaces=true,                 
    numbers=left,                    
    numbersep=5pt,                  
    showspaces=false,                
    showstringspaces=false,
    showtabs=false,                  
    tabsize=4
}

\lstset{style=Cstyle}

% \makebox[\textwidth]{Name:\enspace\hrulefill}
%  
% \vspace{5mm}
%\makebox[\textwidth]{Instructor’s name:\enspace\hrulefill}

\newcommand{\truefalse}
{		\begin{checkboxes}
			 \choice True
			 \choice False. Explanation: 
		\end{checkboxes}
		\vspace{10pt}
}

\newcommand{\answerspace}[1]
{
	\ifprintanswers
	\else
		\setlength{\mylength}{10pt}%
		%\noindent\rule{\mylength}{20pt}
		%\vspace{\fpeval{1.2*(#1)}}
		\vspace{#1\mylength}
	\fi
}
\newlength{\mylength}

\newcommand\pf[1]{\textcolor{red}{PF: #1}}
\newcommand\am[1]{\textcolor{blue}{AM: #1}}

\newcommand{\answer}[1]{\begin{solution}{#1}\end{solution}}
\newcommand{\ra}{$\rightarrow$\xspace}



\usepackage{color}
\usepackage{lastpage}
\usepackage{verbatim}
\usepackage{fancyvrb}
\usepackage{graphicx}
\usepackage{float}
\usepackage{algorithm}
\usepackage{listings}
\usepackage{algpseudocode}
\usepackage{tabularx}
\usepackage{amsmath}

\newcommand{\includecode}[2][c]{\lstinputlisting[caption=, escapechar=, language=#1]{#2}<!---->}
\RecustomVerbatimCommand{\VerbatimInput}{VerbatimInput}{}

\renewcommand{\solutiontitle}{\noindent\textbf{Solution sketch: }}

\definecolor{SolutionColor}{rgb}{0.8,0.9,1}


\cfoot{Page \thepage\ of \pageref{LastPage}}

\rhead{Initials: \rule{20mm}{0.4pt}}



%%%%%%%%%%%%%%%%%%%%%%%%%%%%%%%%%%%%%%%%
\begin{document}
\fvset{numbers=left,numbersep=3pt}
\thispagestyle{empty}

\begin{center}
	{\LARGE CS408 Fall 2025 -- Homework 3 (100 points)\\}

	{\Large Deadline: October 24, Friday\\}

	\vspace{20mm}
\fbox{\fbox{\parbox{5.5in}{\centering
	\begin{itemize}
		\item The quiz is open notes, open book.
		\item No collaboration and no communication of any kind allowed.
		\item Provide complete and clear but concise answers. Justify thoroughly statements where required.
		\item coding questions should be answered in c.
		\item Use basic blocks for all cFG exercises.
		\item The conditions for $p_x$ should be simplified as shown in class.
	\end{itemize}
	}}}
\end{center}

\vspace{13mm}
\begin{center}
	\makebox[0.6\textwidth]{
		\combinedgradetable[h][questions]
	}
\end{center}
%\begin{center}\combinedgradetable[h][questions]\end{center}
\vspace{28mm}

\begin{center}
\makebox[\textwidth]{\textbf{Name (please print)}: \rule{100mm}{0.8pt} }
\end{center}

\newpage
\begin{questions}

\question For the following clause: $p = (a \oplus b) \land (\lnot c)$

\begin{parts}

% correct = 5
% Partial simplification = 3
% Incorrect = 0
\part[5] Provide the final (simplified) $p_a$ condition. Please include simplification steps.
\answerspace{10}

\answer{
	\\
	$p_a = p_{a = \top} \oplus p_{a = \bot}$ \\
	$p_a = \bigl((\top \oplus b) \land \lnot c\bigr) \oplus \bigl((\bot \oplus b) \land \lnot c\bigr)$ \\
	$p_a = \bigl(\lnot b \land \lnot c\bigr) \oplus \bigl(b \land \lnot c\bigr)$ \\
	% factor out common \lnot c
	$p_a = \lnot c \land \bigl(\lnot b \oplus b\bigr)$ \\
	but $\lnot b \oplus b = \top$ (they are complements), hence \\
	\[
	\boxed{\,p_a \;=\; \lnot c\,}
	\]
    \\
}


% correct = 5
% Partial simplification = 3
% Incorrect = 0
\part[5] Provide the final (simplified) $p_b$ condition. Please include simplification steps.
\answerspace{10}
\answer{
	\\
	$p_b = p_{b = \top} \oplus p_{b = \bot}$ \\
	$p_b = \bigl((a \oplus \top) \land \lnot c\bigr) \oplus \bigl((a \oplus \bot) \land \lnot c\bigr)$ \\
	$p_b = \bigl(\lnot a \land \lnot c\bigr) \oplus \bigl(a \land \lnot c\bigr)$ \\
	% factor out common \lnot c
	$p_b = \lnot c \land \bigl(\lnot a \oplus a\bigr)$ \\
	Since $\lnot a \oplus a = \top$, we get \\
	\[
	\boxed{\,p_b \;=\; \lnot c\,}
	\]
}


% correct = 5
% Partial simplification = 3
% Incorrect = 0
\part[5] Provide the final (simplified) $p_c$ condition. Please include simplification steps.
\answerspace{10}
\answer{
	\\
	$p_c = p_{c = \top} \oplus p_{c = \bot}$ \\
	$p_{c=\top} = (a \oplus b)\land \lnot\top = (a \oplus b)\land \bot = \bot$ \\
	$p_{c=\bot} = (a \oplus b)\land \lnot\bot = (a \oplus b)\land \top = (a \oplus b)$ \\
	Therefore
	\[
	p_c = \bot \oplus (a \oplus b) = a \oplus b.
	\]
	Hence
	\[
	\boxed{\,p_c \;=\; a \oplus b\,.}
	\]
}



\begin{table}
    \begin{center}
        \begin{tabular}{||c c c c c||} 
    \hline
        & a & b & c & $(a \oplus b) \land (\lnot c)$ \\ [0.5ex] 
    \hline\hline
     1 & T & T & T & F \\ 
     \hline
     2 & T & T & F & F \\
     \hline
     3 & T & F & T & F \\
     \hline
     4 & T & F & F & T \\
     \hline
     5 & F & T & T & F \\ 
     \hline
     6 & F & T & F & T \\
     \hline
     7 & F & F & T & F \\
     \hline
     8 & F & F & F & F \\ 
     [1ex] 
    \hline
        \end{tabular}
        \caption{Truth table for $(a \oplus b) \land (\lnot c)$.\label{table:truth_xor}}
    \end{center}
\end{table}

% \pf{Please add a question for them to write the truth table, make it worth 4 points. Do it for this and the next predicate.}
\part[4] Write the truth table for the given predicate in this question.
\answerspace{10}
\answer{
	The truth table is Table 2.
}	

% 2 each for the true and false case
\part[4] For the given predicate in this question, determine whether the following test set satisfies Predicate Coverage:
\begin{itemize}
	\item {Test 1: a=True, b=False, c=False}
	\item {Test 2: a=False, b=True, c=True}
\end{itemize}
\answerspace{10}
\answer{
	The truth table is Table 1. \\
	% The given inputs correspond to rows 4 and 5 in the truth table, where row 4 evaluates to True and row 5 evaluates to False. \\
	% To satisfy Predicate Coverage, we need at least one test case that evaluates to True and one that evaluates to False. \\
	Test 1 evaluates the predicate to true, and Test 2 evaluates it to false. 
	Hence, this input satisfies Predicate Coverage.
}

% 0.5 points for each T/F covered for a b c
\part[3] Write the TR for Clause Coverage (CC).
% with respect to $TR_a$, $TR_b$ and $TR_c$.
\answerspace{10}
\answer{
	%The truth table is Table 1.  \\
	TR: a=False, a=True, b=False, b=True, c=False, c=True.
	%For clause a, $TR_a$ is: include any pair (r\_true, r\_false), where r\_true$\in$\{1, 2, 3, 4\} and r\_false$\in$\{5, 6, 7, 8\}. \\
	%For clause b, $TR_b$ is: include any pair (r\_true, r\_false), where r\_true$\in$\{1, 2, 5, 6\} and r\_false$\in$\{3, 4, 7, 8\}. \\
	%For clause c, $TR_c$ is: include any pair (r\_true, r\_false), where r\_true$\in$\{1, 3, 5, 7\} and r\_false$\in$\{2, 4, 6, 8\}. \\
	% So the TR is the union of any one in CC\_a, any one in CC\_b and any one in CC\_c.
	%TR: satisfy $TR_a$, $TR_b$ and $TR_c$ simultaneously.
}


\end{parts}

% \pf{Attreyee, please apply the same type of changes I made to question 1 to question 2. Make sure you apply all the changes.}

\question For the following clause: $p = a \land (\lnot b \rightarrow c)$

\begin{parts}

% correct = 5
% Partial simplification = 3
% Incorrect = 0
\part[5] Provide the final (simplified) $p_a$ condition. Please include simplification steps.
\answerspace{10}
\answer{
	\\
	$p_a = p_{a = \top} \oplus p_{a = \bot}$ \\
	$p_a = \bigl((\top) \land (\lnot b \rightarrow c)\bigr) \oplus \bigl((\bot) \land (\lnot b \rightarrow c)\bigr)$ \\
	$p_a = \bigl(\lnot b \rightarrow c\bigr) \oplus \bot$ \\
	Since $X \oplus \bot = X$, we get \\
	$p_a = \lnot b \rightarrow c$ \\
	Using implication equivalence $\;X\rightarrow Y \equiv \lnot X \lor Y\;$, we have \\
	\[
	\lnot b \rightarrow c \;=\; \lnot(\lnot b)\lor c \;=\; b \lor c.
	\]
	Hence
	\[
	\boxed{\,p_a \;=\; b \lor c\,}
	\]
}

% correct = 5
% Partial simplification = 3
% Incorrect = 0
\part[5] Provide the final (simplified) $p_b$ condition. Please include simplification steps.
\answerspace{10}
\answer{
	\\
	$p_b = p_{b = \top} \oplus p_{b = \bot}$ \\
	$p_{b=\top} = a \land (\lnot \top \rightarrow c) = a \land (\bot \rightarrow c) = a \land \top = a$ \\
	$p_{b=\bot} = a \land (\lnot \bot \rightarrow c) = a \land (\top \rightarrow c) = a \land c$ \\
	Therefore
	\[
	p_b = a \oplus (a \land c).
	\]
	Simplify $a \oplus (a\land c)$:
	\[
	a \oplus (a\land c) = a \land \lnot(a\land c)
	= a \land (\lnot a \lor \lnot c)
	= (a\land\lnot a)\;\lor\;(a\land\lnot c)
	= a\land\lnot c.
	\]
	Hence
	\[
	\boxed{\,p_b \;=\; a \land \lnot c\,}
	\]
}

% correct = 5
% Partial simplification = 3
% Incorrect = 0
\part[5] Provide the final (simplified) $p_c$ condition. Please include simplification steps.
\answerspace{10}
\answer{
	\\
	$p_c = p_{c = \top} \oplus p_{c = \bot}$ \\
	$p_{c=\top} = a \land (\lnot b \rightarrow \top) = a \land \top = a$ \\
	$p_{c=\bot} = a \land (\lnot b \rightarrow \bot) = a \land \lnot(\lnot b) = a \land b$ \\
	Therefore
	\[
	p_c = a \oplus (a \land b).
	\]
	Simplify $a \oplus (a\land b)$ similarly to above:
	\[
	a \oplus (a\land b) = a \land \lnot b.
	\]
	Hence
	\[
	\boxed{\,p_c \;=\; a \land \lnot b\,}
	\]
}

\begin{table}
    \begin{center}
        \begin{tabular}{||c c c c c||} 
    \hline
        & a & b & c & $a \land (\lnot b \rightarrow c)$ \\ [0.5ex] 
    \hline\hline
     1 & T & T & T & T \\ 
     \hline
     2 & T & T & F & T \\
     \hline
     3 & T & F & T & T \\
     \hline
     4 & T & F & F & F \\
     \hline
     5 & F & T & T & F \\ 
     \hline
     6 & F & T & F & F \\
     \hline
     7 & F & F & T & F \\
     \hline
     8 & F & F & F & F \\ 
     [1ex] 
    \hline
        \end{tabular}
        \caption{Truth table for $a \land (\lnot b \rightarrow c)$.\label{table:truth_abc}}
    \end{center}
\end{table}

\part[4] Write the truth table for the given predicate in this question.
\answerspace{10}
\answer{
	The truth table is Table 2.
}	

% 2 each for the true and false case
\part[4] For the given predicate in this question, determine whether the following test set satisfies Predicate Coverage:
\begin{itemize}
	\item {Test 1 : a=True , b=False , c=False}
	\item {Test 2 : a=False , b=True , c=True}
\end{itemize}
\answerspace{10}
\answer{
	The given inputs correspond to rows 4 and 5 in the truth table, which both evaluate to False. \\
	To satisfy Predicate Coverage, we need at least one test case that evaluates to True and one that evaluates to False. \\
	Hence, this input does not satisfy Predicate Coverage.
}

% 0.5 points for each T/F covered for a b c
\part[3] Write the TR for Clause Coverage (CC).
\answerspace{10}
\answer{
	TR: a=False, a=True, b=False, b=True, c=False, c=True.
}
\end{parts}

\question Answer the following questions for the method checkIt() below, which 
prints ``P is true'' if the boolean values of a, b, and c 
satisfy the condition \texttt{(a \&\& b) || (c)}, and ``P isn't true'' otherwise.
\VerbatimInput{hw3/checkIt.java}
\begin{parts}
\part[5] List all predicates in the checkIt method, including predicates that are part of other predicates.
\answerspace{10}
\answer{
	
	$(a \wedge b) \vee c$ \\
	$a \wedge b$ \\
	$a$ \\
	$b$ \\
	$c$ 
}

% One point for each in TR
\part[7] Derive a Active Clause Coverage (ACC) test set for checkIt(). Explain your answer.
\answerspace{10}
\answer{
	% The truth table of the predicate in (a) is: \\
	% \begin{tabular}{l|l|l|l|l}
	% 	& a & b & c & P \\ \hline
	% 1 & F & F & F & F \\ \hline
	% 2 & F & F & T & F \\ \hline
	% 3 & F & T & F & F \\ \hline
	% 4 & F & T & T & F \\ \hline
	% 5 & T & F & F & F \\ \hline
	% 6 & T & F & T & T \\ \hline
	% 7 & T & T & F & T \\ \hline
	% 8 & T & T & T & T
	% \end{tabular}
	The following forms of ACC: GACC, CACC, and RACC were accepted as correct answers.

	Make $a$ major then we need $c=\bot$ and $b=\top$ for $a$ to determine \(p = (a \wedge b) \vee c\).
	\[
	TR_a = \{ (a = \top, b = \top, c = \bot),\; (a = \bot, b = \top, c = \bot)\}
	\]

	Make $b$ major then we need $c=\bot$ and $a=\top$ to determine 
	\(p = (a \wedge b) \vee c\).

	\[
	TR_b = \{ (a = \top, b = \bot, c = \bot),\; (a = \top, b = \top, c = \bot)\}
	\]

	Make $c$ major then we need $(a \wedge b)=\bot$ to determine 
	\(p = (a \wedge b) \vee c\).

	\[
	TR_c = \{ (a = \bot, b = \bot, c = \bot),\; (a = \bot, b = \bot, c = \top)\}
	\]

	This gives us

	\[
	TR = TR_a \cup TR_b \cup TR_c
	\]

	\[
	TR = \left\{
	\begin{aligned}
	& (a = \top, b = \top, c = \bot),\\
	& (a = \bot, b = \top, c = \bot),\\
	& (a = \top, b = \bot, c = \bot),\\
	& (a = \bot, b = \bot, c = \bot),\\
	& (a = \bot, b = \bot, c = \top)
\end{aligned}
\right\}
\]
}

\end{parts}

\if 0 %%%%%%%%%%%%%%%%%%%%%%%%%%%%%%%%%%%%%%%%%
	\question[10]
	\answerspace{30}
	\answer{}

	\question
	\begin{parts}
		\part[6]
		\answerspace{20}
		\answer{}

		\part[2]
		\answerspace{10}
		\answer{}

	\end{parts}
\fi %%%%%%%%%%%%%%%%%%%%%%%%%%%%%%%%%%%%%%%%%

	\newpage
	\begin{center}
		[EMPTY PAGE]
	\end{center}

\end{questions}

\end{document}
